La falta de rendimientos medidos de forma homogénea limita la caracterización productiva a escala parcelaria. La precipitación, los indicadores de vigor vegetal y las capas topográficas aportan evidencia complementaria, aunque no constituyen observaciones directas de cosecha ni identifican por sí solos las causas del rendimiento \parencite{Sharifi2020,Holland2021,Bokusheva2016,Serban2025}.

Con únicamente 138 etiquetas reales, aumentar sin restricciones la complejidad de un predictor puede producir ajustes inestables y aparentes mejoras que no se repiten en parcelas excluidas. Se justifica, por tanto, una arquitectura con componentes regularizados, vecindarios controlados y comparaciones emparejadas. La metodología prioriza el error cuadrático medio de predicción medido mediante LOOCV, pero conserva MAE, sesgo, errores extremos y diagnósticos espaciales para no ocultar defectos del modelo.

La reconstrucción de 59 rendimientos se plantea como una herramienta de apoyo técnico, no como sustitución automática de mediciones de campo o de verificación independiente. La transparencia sobre imputación, selección de modelos e incertidumbre es necesaria para evitar decisiones basadas en una precisión exagerada.
